% Volume III opening, written from the radioactivity source (sec0)
\volumeopening{About This Volume}

A gram of radium sitting on a bench is never quite still. Every second, about \num{3.7e10} of its nuclei
break apart, each one firing out a helium nucleus at about five percent of the speed of light. The
sample stays measurably warmer than its surroundings, year after year, with no fuel and no wire. In
1903, when Pierre Curie and Albert Laborde first measured that heat, it looked like a violation of the
conservation of energy, the principle on which the whole of Volume~I rests. It turned out to be the
first glimpse of an energy store about a million times deeper than chemistry: the binding of the
nucleus.

This volume follows that store from its foundation to its uses. It first establishes what a nucleus is
made of, how big it is, and why its mass is \emph{less} than the mass of its parts. It then asks what
happens when a nucleus is unstable: which particles come out, how much energy they carry, and what rule
governs \emph{when} a given nucleus decays. The answer to the last question is that nothing does; only
the statistics of many nuclei are lawful, and that one fact produces the exponential law of decay. The
same curve of binding energy then explains both the fission reactor and the Sun. The volume ends with
dating rocks and bones, imaging the body, protecting it from radiation, and simulating decay from pure
randomness.

\section*{The Map of the Volume}

The five chapters fall into two parts. Part~A builds the nucleus and its decay; Part~B spends the
energy that Part~A uncovers.

\begin{center}\small
\begin{tabular}{@{}>{\raggedright\arraybackslash}p{0.30\textwidth}>{\raggedright\arraybackslash}p{0.64\textwidth}@{}}\toprule
\textbf{Chapter} & \textbf{What it does}\\\midrule
\multicolumn{2}{@{}l}{\emph{Part A\quad The Nucleus and Its Decay}}\\[1pt]
\chref{ch:nucleus}\quad The Nucleus & size from Rutherford scattering; mass and energy; mass defect,
binding energy and the curve of binding energy per nucleon; the nuclear force\\[2pt]
\chref{ch:decay}\quad Radioactive Decay & $\alpha$, $\beta^-$, $\beta^+$, electron capture and $\gamma$
decay; $Q$-values and how the energy is shared; the valley of stability; decay series\\[2pt]
\chref{ch:decaylaw}\quad The Law of Radioactive Decay & $dN/dt=-\lambda N$ from randomness; half-life,
mean life and activity; branching; decay chains and equilibrium; production at a constant rate\\[3pt]
\multicolumn{2}{@{}l}{\emph{Part B\quad Nuclear Energy and Its Uses}}\\[1pt]
\chref{ch:reactions}\quad Nuclear Reactions: Fission and Fusion & reaction $Q$-values and threshold
energies; fission and the chain reaction; reactors; fusion, the Coulomb barrier and the Sun\\[2pt]
\chref{ch:nucapps}\quad Applications of Radioactivity & radioactive dating; tracers and medical
isotopes; radiation dose and shielding; simulating decay with random numbers\\\bottomrule
\end{tabular}
\end{center}

Two threads tie the chapters together. The binding-energy curve of \chref{ch:nucleus} supplies the
energetics of every decay in \chref{ch:decay} and of every reaction in \chref{ch:reactions}. The decay
law of \chref{ch:decaylaw} underlies every application in \chref{ch:nucapps}, from carbon dating to the
dose received from a medical scan. The chapters can be read in order; a reader short of time can go
from \chref{ch:nucleus} straight to \chref{ch:reactions} for nuclear energy, or from \chref{ch:decay}
straight to \chref{ch:decaylaw} for the decay law, at the cost of a few back-references.

\section*{What This Volume Builds On}

Almost nothing in this volume needs new mechanics. It needs the conservation laws of Volume~I, applied
to particles too small to see.
\begin{itemize}
\item \textbf{Energy} (\chref{ch:energy}): kinetic and potential energy, and the conservation of their
sum. Every $Q$-value is an energy balance, and the distance of closest approach in Rutherford scattering is
a turning point on a potential-energy curve.
\item \textbf{Momentum} (\chref{ch:collisions}): conservation of momentum in collisions and explosions.
It fixes how the energy of a decay is shared between the daughter nucleus and the emitted particle, and
how fast neutrons lose energy in a reactor moderator.
\item \textbf{Angular momentum and inverse-square orbits} (\chref{ch:rotation}, \chref{ch:gravitation},
\chref{ch:kepler}): a charged particle passing a nucleus feels a repulsive inverse-square force, the
mirror image of gravity, so its path is a conic and its angular momentum about the nucleus is conserved.
\item \textbf{Thermal energy} (\chref{ch:kinetic}): the typical kinetic energy of a particle in a gas at
temperature $T$ is of order $k_BT$, which decides whether nuclei in the core of the Sun can approach one
another.
\item \textbf{Exponentials and probability} (\chref{ch:toolkit}): the exponential function, logarithms and
first-order equations of growth and decay (Section~\ref{sec:explog}), and averages over a distribution
(Section~\ref{sec:averages}). The decay law is exactly the equation $dy/dt=-ky$ met there.
\end{itemize}

Some knowledge from Class~12 physics is assumed, and is used without being re-derived.
\begin{itemize}
\item \textbf{Electrostatics:} Coulomb's law; the electric potential energy $q_1q_2/(4\pi\varepsilon_0r)$
of two point charges a distance $r$ apart, where $\varepsilon_0$ is the permittivity of free space; and
the electron-volt as a unit of energy.
\item \textbf{Magnetic force:} a charge $q$ moving with velocity $\vv$ in a magnetic field $\B$ feels the
force $q\vv\times\B$, which bends positive and negative particles in opposite directions.
\item \textbf{Photons:} light of frequency $f$ and wavelength $\lambda_{\text{ph}}$ comes in quanta of
energy $E=hf=hc/\lambda_{\text{ph}}$ and momentum $E/c$, where $h$ is Planck's constant.
\end{itemize}
Three ideas are genuinely new, and each is flagged where it enters. The equivalence of mass and energy,
$E=mc^2$, comes from special relativity and is used as an experimentally established fact
(\chref{ch:nucleus}). The neutrino, a neutral and nearly massless particle, is introduced to save the
conservation of energy and momentum in $\beta$ decay (\chref{ch:decay}). And quantum tunnelling, by which
a particle crosses an energy barrier it classically could not climb, is described but not derived
(\chref{ch:decay} and \chref{ch:reactions}).

\section*{Where This Volume Sits in the Examinations}

The national syllabi have been trimmed in recent cycles, and the table below records what each
examination currently lists. A thorough understanding of every chapter is still recommended: the decay
law is also the first-order rate law of chemical kinetics, and JEE Advanced and the Physics Olympiad
continue to examine it directly.

\begin{center}\small
\begin{tabular}{@{}>{\raggedright\arraybackslash}p{0.52\textwidth}ccc@{}}\toprule
\textbf{Topic} & \textbf{JEE Main} & \textbf{NEET} & \textbf{JEE Advanced}\\\midrule
Nuclear size, mass defect, binding energy, the $\Eb/A$ curve & yes & yes & yes\\
Fission and fusion (energy release) & yes & yes & yes\\
Reactions $\mathrm{X}(a,b)\mathrm{Y}$, $Q$-value, threshold energy & no & no & problem source\\
$\alpha$, $\beta$, $\gamma$ radiation; decay law; half-life, mean life & removed & removed & yes\\
Successive decay, equilibrium, activation & no & no & problem source\\
First-order kinetics (Chemistry) & yes & yes & yes\\
Dating, medical isotopes, dosimetry & enrichment & enrichment & enrichment\\\bottomrule
\end{tabular}
\end{center}

\section*{Learning Outcomes}

After working through this volume, a reader should be able to:
\begin{enumerate}[label=\arabic*.]
\item compute nuclear radii, densities, mass defects and binding energies from tables of atomic masses,
keeping the electron masses correctly accounted for (\chref{ch:nucleus});
\item write balanced equations for $\alpha$, $\beta^-$, $\beta^+$, electron-capture and $\gamma$ decay, and
compute their $Q$-values and the energy shared between the products (\chref{ch:decay});
\item derive and apply $N(t)=N_0e^{-\lambda t}$, the half-life and mean life, activity, branching,
successive decay and production at a constant rate (\chref{ch:decaylaw});
\item compute $Q$-values and threshold energies of nuclear reactions; explain from the $\Eb/A$ curve why
fission and fusion release energy; and estimate reactor fuel use, moderator collisions and stellar power
(\chref{ch:reactions});
\item date samples from parent and daughter amounts, and solve problems on medical dose and shielding
(\chref{ch:nucapps});
\item simulate radioactive decay with random numbers and explain why the exponential law emerges
(\chref{ch:decaylaw} and \chref{ch:nucapps}).
\end{enumerate}

\section*{Units and Data Used Throughout}

Nuclear physics uses a few units chosen to keep the numbers small. Every one is defined here.

\begin{center}\small
\renewcommand{\arraystretch}{1.15}
\begin{tabular}{@{}>{\raggedright\arraybackslash}p{0.24\textwidth}>{\raggedright\arraybackslash}p{0.69\textwidth}@{}}\toprule
\textbf{Unit or constant} & \textbf{Meaning}\\\midrule
electron-volt, \si{\electronvolt} & the kinetic energy an electron gains when accelerated through a
potential difference of \SI{1}{V}: $\SI{1}{\electronvolt}=\SI{1.602e-19}{J}$. Chemical bonds are a few
\si{\electronvolt}.\\
\si{\keV}, \si{\MeV} & \num{e3} and \SI{e6}{\electronvolt}; $\SI{1}{\MeV}=\SI{1.602e-13}{J}$. Nuclear
energies are typically a few \si{\MeV}.\\
femtometre, \si{\fm} & \SI{e-15}{m}, also called a \emph{fermi}. Nuclei are a few femtometres
across.\\
atomic mass unit, \si{\amu} & one-twelfth of the mass of a neutral carbon-12 atom:
$\SI{1}{\amu}=\SI{1.66054e-27}{kg}$, whose rest energy is \SI{931.494}{\MeV}.\\
\si{\MeV}$/c^2$ & a unit of mass: the mass whose rest energy $mc^2$ is \SI{1}{\MeV} (the electron's
mass is $\SI{0.511}{\MeV}/c^2$).\\
becquerel, \si{\becquerel}; curie, \si{\curie} & activity: \SI{1}{\becquerel} is one decay per second;
$\SI{1}{\curie}=\SI{3.7e10}{\becquerel}$, close to the activity of a gram of radium-226
(\chref{ch:decaylaw}).\\
gray, \si{\gray}; sievert, \si{\sievert} & radiation dose: \SI{1}{\gray} is \SI{1}{J} absorbed per
kilogram; the sievert adds a weighting for biological effect (\chref{ch:nucapps}).\\
$h$, $\hbar$ & Planck's constant $h=\SI{6.626e-34}{J.s}$, and $\hbar=h/2\pi$; convenient combinations are
$\hbar c=\SI{197.3}{\MeV.\fm}$ and $hc=\SI{1240}{\electronvolt.\nano\metre}$.\\
$e^2/4\pi\varepsilon_0$ & \SI{1.44}{\MeV.\fm}, so two elementary charges $r$ femtometres apart
have electric potential energy $1.44/r$ in \si{\MeV}.\\\bottomrule
\end{tabular}
\end{center}

The data used throughout are these. The hydrogen \emph{atom} has mass $M(\iso{1}{H})=\SI{1.007825}{\amu}$,
the bare proton $m_p=\SI{1.007276}{\amu}$, the neutron $m_n=\SI{1.008665}{\amu}$ and the electron $m_e=\SI{0.000549}{\amu}$, that is,
$\SI{0.511}{\MeV}/c^2$. Avogadro's number is $N_A=\SI{6.022e23}{mol^{-1}}$, the speed of light
$c=\SI{2.998e8}{m.s^{-1}}$, one year is \SI{3.156e7}{s}, and $R_0=\SI{1.2}{\fm}$.
Atomic masses quoted in the text are from the 2020 Atomic Mass Evaluation, rounded. A glossary of
every technical term appears at the end of the volume.

\section*{Notation Used in This Volume}

Each symbol below has one meaning throughout Volume~III. Two choices differ from many other books and
are worth knowing in advance. The neutron number of a nucleus is always written $A-Z$, never $N$,
because $N$ (or $N(t)$) is reserved for the number of undecayed nuclei in a sample, the central quantity
of \chref{ch:decaylaw}. And $\lambda$, which is the mean free path in Volume~II, is here the decay
constant; a photon's wavelength, where it is needed, is written $\lambda_{\text{ph}}$. A quantity that
returns from an earlier volume keeps its earlier symbol (kinetic energy $K$, Boltzmann's constant $k_B$).

\begin{center}\small
\renewcommand{\arraystretch}{1.12}
\begin{tabular}{@{}>{\raggedright\arraybackslash}p{0.14\textwidth}>{\raggedright\arraybackslash}p{0.31\textwidth}>{\raggedright\arraybackslash}p{0.16\textwidth}>{\raggedright\arraybackslash}p{0.30\textwidth}@{}}\toprule
\textbf{Symbol} & \textbf{Meaning} & \textbf{Symbol} & \textbf{Meaning}\\\midrule
$A$, $Z$ & mass number, atomic (proton) number & $N$, $N(t)$, $N_0$ & undecayed nuclei; at time $t$; at $t=0$\\
$A-Z$ & neutron number & $\lambda$, $\lambda_i$ & decay constant; partial decay constant of channel $i$\\
$\nuc{A}{Z}{X}$, $\iso{A}{X}$ & nuclide notation & $\Thalf$, $T_i$ & half-life; partial half-life of channel $i$\\
$M$, $m$ & atomic mass; particle or nuclear mass & $\tau$ & mean life, $1/\lambda$\\
$m_p$, $m_n$, $m_e$ & proton, neutron and electron masses & $\Act$ & activity (decays per second)\\
$\si{\amu}$ & atomic mass unit (upright) & $\Prod$ & constant production rate\\
$\Delta m$, $\Eb$ & mass defect; binding energy & $f_i$ & branching fraction of channel $i$\\
$S_n$, $S_p$ & neutron, proton separation energy & $\mathrm{X}\to\mathrm{Y}$ & parent to daughter; $N_X$, $\lambda_X$, \dots\\
$Q$ & energy released in a reaction or decay & 1, 2 (subscripts) & two \emph{independent} isotopes in a mixture\\
$K$ & kinetic energy & $T_{\text{phys}}$, $T_{\text{bio}}$, $T_{\text{eff}}$ & physical, biological, effective half-life\\
$R$, $R_0$ & nuclear radius; \SI{1.2}{\fm} & $k$ & neutron multiplication factor\\
$d$, $b$ & distance of closest approach; impact parameter & $\phi$ & fraction of its energy a neutron keeps per collision\\
$U_B$ & Coulomb barrier height & $\mu$, $x_{1/2}$ & attenuation coefficient; half-value layer\\
$\varepsilon_0$ & permittivity of free space & $\Omega$, $\varepsilon$ & detector solid angle; intrinsic efficiency\\
$k_B$, $h$, $\hbar$ & Boltzmann's and Planck's constants; $h/2\pi$ & $D$, $H$, $w_R$ & absorbed dose (\si{\gray}); equivalent dose (\si{\sievert}); weighting factor\\
$\B$, $L_e$ & magnetic field; electron lepton number & $N_A$ & Avogadro's number\\
$p$ & momentum; $p$ also labels the proton & $p$ (simulation) & decay probability per time step in \chref{ch:nucapps}\\\bottomrule
\end{tabular}
\end{center}

Three letters need a word of warning. Lepton number (\chref{ch:decay}) is written $L_e$, with its
subscript, because plain $L$ is angular momentum in Volume~I; where angular momentum appears in this
volume (Rutherford scattering in \chref{ch:nucleus}) it is written out in full as $mvb$. And $k$, the spring constant in
Volume~I, is here only the neutron multiplication factor of a reactor (\chref{ch:reactions}). The letter
$p$ is a momentum, and in reactions such as $(\alpha,p)$ it names the proton; in the simulation of
\chref{ch:nucapps}, where no momentum appears, it is the probability that a nucleus decays in one time
step.
